\section{Bernstein approximation, Hausdorff moments and exchangeability}
\label{sec:annex-p15-feller-bernstein-hausdorff}

\begin{annexitem}{annex:p15.feller-concentration}{thm:p15-feller-concentration}{The concentration estimate with its uniform domain}
\begin{proof}
Let \(M=\sup_{\mathbb R}|u|<\infty\).
For fixed \(\theta\) and \(\epsilon>0\), continuity gives
\(\delta>0\) such that \(|u(x)-u(\theta)|<\epsilon\) when
\(|x-\theta|<\delta\). Split the integral over this neighborhood
and its complement. Chebyshev's inequality gives
\[
 \left|\int u\,d\nu_{n,\theta}-u(\theta)\right|
       \leq\epsilon+2M\,v_n(\theta)/\delta^2.
\]
This tends to at most \(\epsilon\), proving the pointwise assertion.

For \(\theta\) in a compact interval \(K\), restrict \(\delta\leq1\).
The function is uniformly continuous on the compact set
\(K+[-1,1]\), so the same \(\delta\) works for all \(\theta\in K\).
Uniformly vanishing variances give the uniform conclusion.
If \(u\) is uniformly continuous on the whole line, the same choice
works for any parameter set, so the identical estimate proves the
last assertion. No interchange of an unbounded test integral was used.
\end{proof}
\end{annexitem}

\begin{annexitem}{annex:p15.feller-kernel-examples}{thm:p15-feller-kernel-examples}{The four source families and the gamma variance}
\begin{proof}
For a binomial \(J\) with parameters \(n,\theta\), the scaled variable
\(J/n\) has mean \(\theta\) and variance \(\theta(1-\theta)/n\leq1/(4n)\).
Its finite expectation is the first displayed sum.
A continuous function on \([0,1]\) extends boundedly and continuously
to the line by constant values outside that interval.
Apply the preceding lemma uniformly on \([0,1]\).

For a Poisson variable \(J\) of mean \(n\theta\), the variable \(J/n\)
has mean \(\theta\) and variance \(\theta/n\).
These identities hold for every real \(n>0\); only the Poisson count
\(J\) is an integer. The displayed sum is its expectation, including
\(\theta=0\), where it is \(u(0)\).
A bounded continuous function on the nonnegative half-line extends
constantly to negative arguments. On \(0\leq\theta\leq T\), the
variance is bounded by \(T/n\), proving the assertion.

For the gamma density, the displayed integrand is the density with
shape \(n\) and rate \(n/\theta\).
The already proved gamma formulas give mean
\(n/(n/\theta)=\theta\) and variance
\(n/(n/\theta)^2=\theta^2/n\).
They apply for every positive real \(n\).
At \(\theta=0\) the separate point mass has variance zero.
On \(0\leq\theta\leq T\) the bound is \(T^2/n\), so the same lemma
gives uniform convergence. The integer factorial identity follows
from the gamma recursion.

Finally an average of \(n\) independent observations with common
mean \(\theta\) and variance \(\sigma^2<\infty\) has mean \(\theta\)
and variance \(\sigma^2/n\). The lemma gives \(E u(M_n)\to u(\theta)\).
This is the stated asymptotic unbiasedness. It neither asserts exact
finite-sample unbiasedness nor removes the test-function hypotheses.
\end{proof}
\end{annexitem}

\begin{annexitem}{annex:p15.feller-differences}{thm:p15-feller-differences}{The source reciprocity coefficient by coefficient}
\begin{proof}
The finite-difference formula follows by induction on \(r\).
Applying \(\Delta\) to the formula at \(r\), shifting one sum and
combining adjacent binomial coefficients gives the formula at \(r+1\).
The case \(r=0\) is the identity.
Apply this to the sequence \(a_j=u(x+jh)\) and divide by \(h^r\)
to obtain the normalized function formula.

Expand the left side of the reciprocity formula with that difference
identity. The coefficient of \(a_{i+j}\) is
\[
 \sum_{r=j}^N\binom Nr\binom rj(-1)^{r-j}c_r
   =\binom Nj\sum_{\ell=0}^{N-j}
                    \binom{N-j}{\ell}(-1)^\ell c_{j+\ell}
   =\binom Nj(-1)^{N-j}\Delta^{N-j}c_j.
\]
The middle equality uses
\(\binom Nr\binom rj=\binom Nj\binom{N-j}{r-j}\).
This proves the finite identity without any infinite rearrangement.

Take \(a_i=1\) for every \(i\).
Only its zeroth difference survives on the left, giving \(c_0\).
Replace \(c_j\) by \(c_{k+j}\) to obtain the full shifted inversion.
For \(c_r=\theta^r\), direct differencing gives
\(\Delta^{N-j}c_j=\theta^j(\theta-1)^{N-j}\).
Substitution proves the binomial specialization, including the
endpoints because exponent-zero powers are defined as one.
\end{proof}
\end{annexitem}

\begin{annexitem}{annex:p15.feller-bernstein}{thm:p15-feller-bernstein}{Dual representation and the explicit approximation}
\begin{proof}
In the binomial reciprocity identity take \(i=0\),
\(a_j=u(jh)\), \(h=1/n\). Its sequence difference is
\(\Delta^r a_0=h^r\Delta_h^r u(0)\).
The resulting two sides are exactly the two displayed forms of
\(B_{n,u}\).

For \(J\sim\operatorname{Binomial}(n,\theta)\), split
\(E|u(J/n)-u(\theta)|\) at \(|J/n-\theta|\leq\delta\).
The inside contribution is at most \(\omega_u(\delta)\);
outside it is at most \(2\|u\|_\infty\) times the tail probability.
The variance bound and Chebyshev give the stated estimate.
For instance, \(\delta=n^{-1/4}\) makes its second term tend to
zero and its first term tend to zero by uniform continuity.
If \(u\) is \(L\)-Lipschitz, Cauchy--Schwarz gives instead
\[
 |B_{n,u}(\theta)-u(\theta)|
     \leq L E|J/n-\theta|
     \leq L\sqrt{\theta(1-\theta)/n}\leq L/(2\sqrt n).
\]
For \([a,b]\), apply the approximation to
\(v(t)=u(a+(b-a)t)\), then substitute \(t=(x-a)/(b-a)\).
This is a polynomial in \(x\) and has the same uniform error.
\end{proof}
\end{annexitem}

\begin{annexitem}{annex:p15.feller-positive-series}{thm:p15-feller-positive-series}{Diagonal coefficient selection and the endpoint mass}
\begin{proof}
All coefficients lie in \([0,M]\).
Successive subsequence selection and a diagonal sequence give a
subsequence along which every coefficient converges, say to \(a_r\geq0\).
Every finite partial sum of these limits is at most \(M\), hence
\(\sum_r a_r\leq M\).
For \(0\leq x\leq t<1\), the tail beyond degree \(R\) of any
one of the series is bounded by \(Mt^{R+1}\); the same bound
holds for the limiting coefficients.
Thus finite coefficient convergence followed by this uniform tail
bound gives
\(u(x)=\sum_r a_rx^r\), uniformly on each \([0,t]\).

For uniqueness, if two such series agreed but had a first different
coefficient at degree \(r\), subtract and divide by \(x^r\).
As \(x\downarrow0\), the tails tend to zero, forcing that coefficient
difference to be zero, a contradiction.
Every subsequence has a further diagonal subsequence with these same
coefficients. If a full coefficient sequence failed to converge, a
subsequence staying a fixed distance away would contradict that fact.
Thus \(a_{n,r}\to a_r\) for every \(r\).

Nonnegative terms give
\(\lim_{x\uparrow1}\sum_r a_rx^r=\sum_r a_r\), by monotone convergence
for the counting measure on the indices.
A continuous extension therefore has value \(\sum_r a_r\) at one.
For probability generating functions this recovers total mass one
under the stated endpoint hypothesis.
For \(P_n(x)=x^n\), the limit on \([0,1)\) is zero while the
finite generating functions equal one at \(1\); continuity of
the proposed full endpoint limit fails. This is the lost-mass case.
\end{proof}
\end{annexitem}

\begin{annexitem}{annex:p15.feller-absolute-monotone}{thm:p15-feller-absolute-monotone}{The exact source grid condition without the missing endpoint step}
\begin{proof}
Suppose first \(u(x)=\sum_r p_rx^r\), with \(p_r\geq0\) and
finite sum at each \(x<1\). For \(0<a<b<1\), convergence at \(b\)
bounds every differentiated series uniformly on \([0,a]\):
for a fixed derivative order \(k\), the factor
\((r)_k(a/b)^{r-k}\) is bounded in \(r\), where
\((r)_k=r(r-1)\cdots(r-k+1)\).
Comparison with \(\sum_rp_rb^r\) proves uniform convergence.
Integration of the uniformly convergent derivative series, starting
from the polynomial partial sums, justifies termwise differentiation.
Every derivative is nonnegative on \((0,1)\), including order zero.

Conversely assume all derivatives there are nonnegative.
For \(x>0\) and \(x+kh<1\), repeated use of the fundamental theorem
of calculus gives
\[
 \Delta_h^ku(x)
  =h^{-k}\int_{[0,h]^k}u^{(k)}(x+t_1+\cdots+t_k)\,
                                  dt_1\cdots dt_k\geq0
\]
for \(k\geq1\). Fix \(h=1/n\), \(k<n\), then let \(x\downarrow0\).
Continuity of \(u\) at each finite grid point makes the finite
difference converge to \(\Delta_{1/n}^ku(0)\), which is nonnegative.
The \(k=0\) conclusion follows from \(u(0)=\lim_{x\downarrow0}u(x)\geq0\).
No derivative at zero was assumed.

It remains to deduce the positive series from the grid hypothesis.
Fix \(0<a<1\), let \(N\to\infty\) be integers, and put
\(m_N=\lfloor aN\rfloor\).
For all sufficiently large \(N\), \(1\leq m_N<N\).
Use the polynomial
\[
 P_N(x)=\sum_{j=0}^{m_N}u(j/N)\binom{m_N}j x^j(1-x)^{m_N-j}.
\]
The reciprocity formula gives its power coefficients as
\[
 P_N(x)=\sum_{r=0}^{m_N}\binom{m_N}r N^{-r}
                                      \Delta_{1/N}^ru(0)\,x^r.
\]
Every coefficient is nonnegative: only \(r\leq m_N<N\) is used.
Moreover their sum is \(P_N(1)=u(m_N/N)\), bounded by
\(\max_{[0,a]}|u|\).
For a binomial \(J\) of parameters \(m_N,x\), the variable \(J/N\)
has mean \((m_N/N)x\) and variance at most \(m_N/(4N^2)\).
Uniform continuity on \([0,a]\) and the concentration estimate give
\(P_N(x)-u((m_N/N)x)\to0\) uniformly for \(0\leq x\leq1\).
Also \(u((m_N/N)x)\to u(ax)\) uniformly, since
\(|m_N/N-a|\leq1/N\).
The positive-series lemma therefore gives
\[
                         u(ax)=\sum_{r\geq0}b_r^{(a)}x^r,
                         \qquad b_r^{(a)}\geq0.
\]
This works even if \(u(a)=0\); no division by \(u(a)\) is made.

Equivalently, \(u(y)=\sum_r b_r^{(a)}a^{-r}y^r\) for \(0\leq y<a\).
Rescale a smaller closed interval inside the common domain of two
such expansions to \([0,1]\). Both coefficient sums are then finite,
so the uniqueness argument in the positive-series lemma applies.
It makes the coefficients agree for different \(a\).
Take a sequence \(a\uparrow1\) to obtain one nonnegative coefficient
sequence representing \(u\) throughout \([0,1)\).
This proves the original \(k<n\) criterion rather than assuming
the additional \(k=n\) condition or silently rescaling its grid.
Finally the sum of the coefficients is \(u(1-)\) by monotone
convergence, so total mass one is exactly the generating-function
condition. The zero and unbounded cases are included.
\end{proof}
\end{annexitem}


\begin{annexitem}{annex:p15.feller-hausdorff-weights}{thm:p15-feller-hausdorff-weights}{Positive grid weights and their polynomial limits}
\begin{proof}
For moments of a probability on \([0,1]\), finite binomial expansion gives
\[
 (-1)^r\Delta^r c_k
   =\int\sum_{j=0}^r\binom rj(-1)^j x^{k+j}\,\mu(dx)
   =\int x^k(1-x)^r\,\mu(dx).
\]
This is nonnegative, including \(k=0\) and \(r=0\) under the power
convention. Its exponent is \(r\), with no subtraction of \(k\).

For an arbitrary completely monotone sequence, the displayed weights
are nonnegative by their definition.
The inversion identity in the finite-difference theorem gives
\[
                 \sum_{j=0}^nw_{n,j}=c_0=1.
\]
Thus \(\mu_n\) is a probability, without an assumed original measure.
Apply reciprocity with \(a_j=u(jh)\), \(h=1/n\), to get
\[
 \sum_{j=0}^n u(jh)\binom nj(-1)^{n-j}\Delta^{n-j}c_j
               =\sum_{r=0}^n\binom nr c_r h^r\Delta_h^ru(0).
\]
This is the exact expectation formula.

For a polynomial of degree \(d\), every difference of order greater
than \(d\) is zero: taking a first difference lowers the degree
of a nonconstant polynomial, and constants have zero first difference.
For fixed \(r\leq d\), the repeated integral formula for function
differences gives \(\Delta_h^ru(0)\to u^{(r)}(0)\) as \(h\downarrow0\);
the zeroth case is already equality. Also
\[
                   \binom nr n^{-r}
                     =\frac1{r!}\prod_{\ell=0}^{r-1}(1-\ell/n)
                          \longrightarrow\frac1{r!}.
\]
For \(n\geq d\) there are only \(d+1\) nonzero terms, so their limits
give the claimed polynomial formula.
Taking \(u(x)=x^k\) gives \(c_k\), since its only nonzero derivative
at zero among these terms is the \(k\)th, with value \(k!\).
\end{proof}
\end{annexitem}

\begin{annexitem}{annex:p15.feller-ramp-measure}{thm:p15-feller-ramp-measure}{Constructing the CDF before representing the functional}
\begin{proof}
For continuous \(u\) and \(\epsilon>0\), choose a polynomial \(p\)
with \(\|u-p\|_\infty<\epsilon\), using Bernstein approximation.
Since \(\mu_n\) are probabilities,
\[
                     |E_nu-E_mu|
                        \leq2\epsilon+|E_np-E_mp|.
\]
The polynomial expectations converge by the preceding theorem,
so the sequence \(E_nu\) is Cauchy and has a limit \(L(u)\).
Linearity, positivity, \(L(1)=1\) and the uniform norm bound
follow by taking limits of the corresponding finite-measure identities.

The ramps are continuous on \([0,1]\) and lie in \([0,1]\).
They increase with \(t\), and decrease as \(\epsilon\downarrow0\).
Thus \(F(t)=\inf_{\epsilon>0}L(u_{t,\epsilon})\) exists in \([0,1]\)
and is nondecreasing. If \(t<0\), choose \(\epsilon<-t\);
the ramp is identically zero, so \(F(t)=0\).
If \(t\geq1\), every ramp is identically one, so \(F(t)=1\).

For fixed \(\epsilon>0\), the clipped linear ramps satisfy
\(\|u_{s,\epsilon}-u_{t,\epsilon}\|_\infty\leq|s-t|/\epsilon\).
For \(s>t\),
\[
           F(t)\leq F(s)\leq U_\epsilon(s)
                      \leq U_\epsilon(t)+(s-t)/\epsilon.
\]
Let \(s\downarrow t\), and then take the infimum over \(\epsilon>0\).
The bounds give \(F(t+)=F(t)\).
This proves right continuity directly, rather than assuming that a
limit of continuous functions must be continuous.

The Stieltjes construction \cref{thm:p15-stieltjes-construction}
now gives the measure with interval masses \(F(b)-F(a)\).
Continuity of a measure on increasing intervals and the established
limits at infinity give total mass one and cumulative function \(F\).
It is supported on \([0,1]\), because \(F(t)=0\) for \(t<0\)
and \(F(1)=1\). A possible mass at zero is \(F(0)\), not discarded.

For any real \(t\) and \(0<\epsilon<\delta\), pointwise on \([0,1]\),
\[
                  u_{t-\delta,\epsilon}
                          \leq\mathbf1_{[0,t]}\leq u_{t,\epsilon}.
\]
Here the indicator is understood as \(\mathbf1_{\{x\leq t\}}\)
on the interval, including its empty and whole-interval cases.
Integrating against \(\mu_n\) gives
\[
                  E_nu_{t-\delta,\epsilon}
                         \leq F_n(t)\leq E_nu_{t,\epsilon}.
\]
Take limits inferior and superior in \(n\), then let
\(\epsilon\downarrow0\). This yields
\[
            F(t-\delta)\leq\liminf_nF_n(t)
                     \leq\limsup_nF_n(t)\leq F(t).
\]
If \(F\) is continuous at \(t\), let \(\delta\downarrow0\).
The two bounds coincide, proving the required CDF convergence.
\end{proof}
\end{annexitem}

\begin{annexitem}{annex:p15.feller-hausdorff}{thm:p15-feller-hausdorff}{Moments, uniqueness and the full inversion formula}
\begin{proof}
Let \(\mu\) and its CDF \(F\) be the measure just constructed.
For any probability \(\eta\) supported on \([0,1]\), with CDF \(G\),
and integer \(r\geq1\), the identity
\[
                  x^r=\int_0^1r t^{r-1}\mathbf1_{\{t<x\}}\,dt
                                      \quad(0\leq x\leq1)
\]
and Tonelli's theorem give
\[
               \int x^r\,\eta(dx)
                  =r\int_0^1 t^{r-1}(1-G(t))\,dt
                  =1-r\int_0^1t^{r-1}G(t)\,dt.
\]
This proves the exact source integration-by-parts identity for every
\(r\geq1\), including both endpoint atoms.

Apply it first to \(\mu_n\).
The CDFs converge to \(F\) at every continuity point.
A CDF has at most countably many discontinuities, by
\cref{thm:p15-cdf-laws}, so convergence holds for Lebesgue-almost
every \(t\). The dominating function \(r t^{r-1}\) is integrable.
Dominated convergence therefore gives
\[
                     c_r=\lim_n E_n(x^r)
                                  =\int x^r\,\mu(dx).
\]
The zeroth moment is its total mass one.
This proves existence from complete monotonicity.

The polynomial limit formula agrees with integration against \(\mu\).
For continuous \(u\), uniformly approximate it by a polynomial \(p\).
The norm bound for \(L\) and for probability integration gives
\[
                  |L(u)-\int u\,d\mu|\leq2\|u-p\|_\infty.
\]
Thus \(L(u)=\int u\,d\mu\).
If another probability \(\nu\) on \([0,1]\) has the same moments,
the same polynomial approximation gives \(L(u)=\int u\,d\nu\).
Apply this to each ramp and let \(\epsilon\downarrow0\).
Bounded convergence gives the same CDF \(F(t)\) for \(\nu\)
at every real \(t\). CDF uniqueness proves \(\nu=\mu\).

The explicit sum in the statement is exactly \(F_n(t)\), since
\(j/n\leq t\) is equivalent to \(j\leq nt\).
The preceding ramp theorem therefore proves its convergence at
every continuity point, which is Feller's inversion formula.
Moreover, from the already established moments and difference identity,
\[
            c_n=\int x^n\,d\mu\longrightarrow\mu(\{1\}),\qquad
            (-1)^n\Delta^nc_0=\int(1-x)^n\,d\mu
                                      \longrightarrow\mu(\{0\})
\]
by bounded convergence. These endpoint formulas explain explicitly
how both endpoint masses are retained.

For finite \(c_0>0\), normalize the sequence by \(c_0\), construct
its probability, then multiply that measure by \(c_0\).
Linearity of the differences and uniqueness give the finite-measure
version. If \(c_0=0\), complete monotonicity gives
\(0\leq c_{k+1}\leq c_k\); induction forces every term to be zero.
The only nonnegative finite measure of mass zero is the zero measure.
\end{proof}
\end{annexitem}

\begin{annexitem}{annex:p15.feller-moment-calculations}{thm:p15-feller-moment-calculations}{Computed weights and the finite-prefix limitation}
\begin{proof}
Uniform probability on \([0,1]\) has moment \(1/(k+1)\).
The difference integral formula gives
\[
 w_{n,j}=\binom nj\int_0^1x^j(1-x)^{n-j}\,dx
         =\binom nj B(j+1,n-j+1)=\frac1{n+1},
\]
using the proved beta-gamma identity and factorial values.
Thus the weights and the claimed law are explicit.

For \((1-a)\delta_0+a\delta_1\), all positive-order moments are \(a\).
At zero the weight integrand contributes only when \(j=0\);
at one it contributes only when \(j=n\).
Hence the weights are \(1-a\) and \(a\) at the two endpoints and
zero elsewhere, for every \(n\geq1\).

For \(\delta_{t_0}\), the weights are
\(\binom njt_0^j(1-t_0)^{n-j}\), so \(\mu_n\) is the law of
a binomial count divided by \(n\).
For \(0<t_0<1\), apply the already proved Bernoulli central limit
theorem to the centered count divided by
\(\sqrt{nt_0(1-t_0)}\).
The threshold for its event \(J\leq nt_0\) is
\[
              \frac{\lfloor nt_0\rfloor-nt_0}
                         {\sqrt{nt_0(1-t_0)}}\longrightarrow0.
\]
For each \(\epsilon>0\) it eventually lies between
\(-\epsilon\) and \(\epsilon\).
The central limit bounds at these fixed points, followed by
\(\epsilon\downarrow0\), give limiting probability \(\Phi(0)=1/2\).
This supplementary calculation is not a premise of the Hausdorff
construction. The target CDF at its atom is one, as claimed.

For \(a<b\), push a law on \([a,b]\) through
\(y=(x-a)/(b-a)\).
The finite binomial expansion gives exactly the displayed \(d_k\),
so necessity follows from Hausdorff.
Conversely construct the law of \(Y\in[0,1]\) with moments \(d_k\)
and put \(X=a+(b-a)Y\).
Substitution of the formula for \(d_j\) into the binomial expression
for \(EX^k\) gives coefficient of \(c_\ell\)
\[
 \binom k\ell\sum_{r=0}^{k-\ell}\binom{k-\ell}r
                         a^{k-\ell-r}(-a)^r .
\]
It is zero for \(\ell<k\) and one for \(\ell=k\), by the binomial
identity. Thus \(EX^k=c_k\).
For \(a=b\), the interval supports only \(\delta_a\), giving the
separate criterion without division by \(b-a\).

For the final prefix, the order-zero terms are \(1,1/2,1/10\);
the available negative first differences are \(1/2,2/5\);
the second difference is \(1/10\). All are nonnegative.
Yet any representing probability would have variance
\(1/10-(1/2)^2=-3/20\), impossible by nonnegativity of a squared
centered variable. Thus these finite difference tests cannot by
themselves establish a representing measure.
\end{proof}
\end{annexitem}


\begin{annexitem}{annex:p15.feller-lognormal-moments}{thm:p15-feller-lognormal-moments}{The source's shifted odd integral and moment-growth condition}
\begin{proof}
With \(t=\log x\), \(dx/x=dt\), the base density integrates to the
normalized standard Gaussian integral. Completing the square gives
\[
 \int_0^\infty x^k f(x)\,dx
    =\frac1{\sqrt{2\pi}}\int_{\mathbb R}e^{kt-t^2/2}\,dt
    =e^{k^2/2}\quad(k\geq0).
\]
For \(|a|\leq1\), \(f_a\geq0\) and \(f_a\leq2f\).
All nonnegative integer moment integrals are therefore finite.
The difference from the base moment is proportional to
\[
 \frac1{\sqrt{2\pi}}\int_{\mathbb R}
             e^{kt-t^2/2}\sin(2\pi t)\,dt.
\]
Set \(t=y+k\). Since \(k\) is an integer,
\(\sin(2\pi(y+k))=\sin(2\pi y)\), while the Gaussian factor
becomes \(e^{k^2/2}e^{-y^2/2}\).
The remaining integrand is odd and absolutely integrable, so its
integral is zero. This proves every moment identity, including
normalization at \(k=0\), by Feller's exact cancellation route.

For distinct \(a,b\), the difference of densities has fixed nonzero
sign on \(1<x<e^{1/2}\), where \(f(x)>0\) and
\(\sin(2\pi\log x)>0\). Its integral there is nonzero.
Thus the probability laws really differ despite identical moments.

For the sufficient moment bound let the displayed series have finite
sum \(C\), and put \(\rho=\sqrt s>0\).
If all positive-order even moments are positive, its terms imply
\[
                m_{2n}\leq C(2n)!\rho^{-2n}
                           \leq C(2n)^{2n}\rho^{-2n}.
\]
Hence
\(m_{2n}^{-1/(2n)}\geq \rho/(2nC^{1/(2n)})\).
The Carleman series diverges by comparison with the harmonic series,
so \cref{thm:p15-criteria-carleman} proves determinacy.
The same estimate also meets the growth hypothesis of the independently
proved analytic criterion \cref{thm:p15-criteria-analytic-determinacy}.
If an even moment of positive order vanishes, the nonnegative
variable \(X^{2n}\) is zero almost surely and the law is \(\delta_0\).
This verifies the source's sufficient condition without claiming its
later transform-based source proof has been reproduced here.
\end{proof}
\end{annexitem}

\begin{annexitem}{annex:p15.feller-binary-mixture}{thm:p15-feller-binary-mixture}{The complete-monotonicity route from all-success probabilities}
\begin{proof}
For \(k<n\), partition the first \(n-1\) word event according to its
last unspecified coordinate. The extension with a zero has probability
\(p_{k,n}\). The extension with a one has probability \(p_{k+1,n}\)
by moving that one next to the first \(k\) ones, using exchangeability.
Therefore
\[
                             p_{k,n}=p_{k,n-1}-p_{k+1,n}.
\]
The base values are \(p_{k,k}=c_k\).
Induction on \(r=n-k\) now gives
\[
 p_{k,k+r}
    =(-1)^{r-1}\Delta^{r-1}c_k
       -(-1)^{r-1}\Delta^{r-1}c_{k+1}
    =(-1)^r\Delta^r c_k.
\]
The \(r=0\) case is the base equality.
Because the sequence is infinite, every \(k,r\geq0\) occurs, and
these probabilities are nonnegative. With \(c_0=1\), Hausdorff's
theorem gives a unique \(\lambda\) whose moments are \(c_k\).
Its difference-integral identity gives
\[
                 (-1)^{n-k}\Delta^{n-k}c_k
                         =\int\theta^k(1-\theta)^{n-k}\,\lambda(d\theta).
\]
Every word of length \(n\) and count \(k\) has that probability.
The \(\binom nk\) disjoint words give the count formula.

Conversely choose a parameter \(\Theta\) with law \(\lambda\) and
a mutually independent uniform sequence independent of \(\Theta\).
Set \(X_j=\mathbf1_{\{U_j\leq\Theta\}}\).
Finite product integration gives the displayed word probabilities,
including parameters zero and one. They depend only on the count,
so every finite vector is exchangeable.
Its all-success probabilities are exactly the moments of \(\lambda\),
and Hausdorff uniqueness proves uniqueness of this mixing law.
The construction and proof have used the source's moment route rather
than the general reverse conditional-product theorem.
\end{proof}
\end{annexitem}

\begin{annexitem}{annex:p15.feller-polya-moments}{thm:p15-feller-polya-moments}{The gamma ratio identifies the source mixing law}
\begin{proof}
On the event that every first \(n\) draw is black, the successive
black weights are \(b,b+c,\ldots,b+(n-1)c\), and the total weights
are \(b+r,b+r+c,\ldots,b+r+(n-1)c\).
Their conditional probability product is the first displayed ratio.
Divide its factors by \(c>0\) and apply
\(\Gamma(x+n)/\Gamma(x)=x(x+1)\cdots(x+n-1)\)
to obtain the gamma expression.

The same numerator-permutation argument as in the direct urn proof
shows that its draw sequence is exchangeable.
The Feller binary theorem therefore supplies a mixing probability
whose moments are these \(c_n\).
For \(\alpha=b/c,\beta=r/c\), the normalized beta law has moments
\[
                   \frac{B(\alpha+n,\beta)}{B(\alpha,\beta)}
                    =\frac{\Gamma(\alpha+n)\Gamma(\alpha+\beta)}
                            {\Gamma(\alpha+\beta+n)\Gamma(\alpha)}.
\]
They equal \(c_n\), so Hausdorff uniqueness identifies the mixing law.
Its mixed moment is
\(B(\alpha+k,\beta+n-k)/B(\alpha,\beta)\).
This is the probability of each specified word of count \(k\);
multiplication by \(\binom nk\) gives the count probability.
These agree algebraically with \cref{thm:p15-polya-word-law},
after renaming the first color and its initial weight.
This retains the source's identification through moments.
\end{proof}
\end{annexitem}

\begin{annexitem}{annex:p15.feller-finite-alphabet}{thm:p15-feller-finite-alphabet}{The explicit simplex and directing-law bridge}
\begin{proof}
The finite discrete alphabet is standard Borel, so
\cref{thm:p15-definetti-kernel} supplies the directing probability
\(K\). Map it to the vector
\(p=(K(\{1\}),\ldots,K(\{d\}))\).
Every such vector belongs to \(\Delta_d\), and the map is measurable.
Conversely \(p\mapsto\sum_i p_i\delta_i\) is a measurable inverse:
evaluation at any subset of the finite alphabet is the finite sum
of its corresponding coordinates.
Thus the unique directing-law distribution becomes a unique probability
\(\lambda\) on \(\Delta_d\).

For a specified word, conditional independent identical distribution
makes its probability the product of the corresponding coordinate
masses. Grouping equal letters gives \(\prod_i p_i^{n_i}\);
integration gives the displayed formula.
Choosing positions for the letters successively gives
\[
             \binom n{n_1}\binom{n-n_1}{n_2}\cdots
                            =\frac{n!}{\prod_i n_i!}
\]
possible words of those counts.
They are disjoint and have the same probability, proving the count law.
The empirical frequencies of the finitely many singletons converge
simultaneously to the coordinates of \(K\), by
\cref{thm:p15-directing-law-uniqueness}.
This proves the asserted limit vector and its distribution.
For \(d=1\), the simplex is a singleton and the same formulas give
the deterministic one-letter sequence.
This is the stated host bridge for Feller's finite-alphabet generalization.
\end{proof}
\end{annexitem}

\begin{annexitem}{annex:p15.feller-finite-exchangeability}{thm:p15-feller-finite-exchangeability}{The occupancy and full residual obstructions}
\begin{proof}
The six equally likely allocations are
\[
 (2,0,0),\ (0,2,0),\ (0,0,2),\
 (1,1,0),\ (1,0,1),\ (0,1,1).
\]
Their occupancy indicators are the corresponding three unit vectors
and the three vectors with two ones. Permuting cells permutes this
list, so it is exchangeable.
Exactly three outcomes occupy the first cell, one occupies both
the first two, and none occupies all three.
This gives \(c_1=1/2,c_2=1/6,c_3=0\).

A completely monotone infinite continuation would be nonnegative
and nonincreasing, so every later \(c_j\) would be zero.
Then
\[
               \Delta^4c_1=c_5-4c_4+6c_3-4c_2+c_1
                                         =-\frac16<0,
\]
contradicting complete monotonicity.
An infinite exchangeable extension would create precisely such a
continuation by the Feller binary theorem.
An identical-independent binary mixture for these three coordinates
would likewise give an infinite moment continuation
\(c_j=\int\theta^j\,d\lambda\).
Thus neither exists.

For the residual example write \(\overline Z=n^{-1}\sum_i Z_i\).
The map to the residual vector commutes with permutations, so the
full residual vector, and its first \(n-1\) coordinates, are exchangeable.
Suppose, for a contradiction, that their law is
\(\int\mu^{\otimes(n-1)}\,\lambda(d\mu)\) for a mixing probability
on the measurable space of real probability laws.
Put
\[
 \phi(t)=E e^{itZ_1},\qquad c(t)=|\phi(t)|^2,\qquad
 q_\mu(t)=\int e^{itx}\,\mu(dx).
\]
These transforms are continuous, \(c(0)=1\), and \(|q_\mu(t)|\leq1\).
All the integrals in the argument below are bounded, so no moment
hypothesis is used.
Since \(Y_1-Y_2=Z_1-Z_2\),
\[
                         \int|q_\mu(t)|^2\,\lambda(d\mu)=c(t).
\]
There is also the exact cancellation
\[
                  Y_1+\cdots+Y_{n-2}+2Y_{n-1}=Z_{n-1}-Z_n.
\]
Independence of the original \(Z_i\) and the assumed product-mixture
law therefore give
\[
                \int q_\mu(t)^{\,n-2}q_\mu(2t)\,\lambda(d\mu)=c(t).
                                                        \tag{*}
\]

We first record the elementary transform fact used below.
If independent identically distributed real \(V,V'\) satisfy
\(E e^{it(V-V')}=1\) on a neighborhood of zero, then \(V=V'\)
almost surely. Indeed, for countably many rational \(t\) in that
neighborhood, the nonnegative variables \(1-\cos(t(V-V'))\)
have expectation zero. On their common full-measure event a nonzero
difference \(d\) is impossible: choose a nonzero rational \(t\) with
\(0<|td|<2\pi\). Thus the difference is zero.
An i.i.d. pair is equal almost surely only when its common law is
a point mass. To see this, every Borel set \(B\) would satisfy
\(P(V\in B)P(V'\notin B)=0\); otherwise the pair differs with
positive probability. Hence every rational half-line has probability
zero or one. Its CDF then concentrates at a single real threshold,
using its endpoint limits and right continuity.
In particular a characteristic function whose modulus is one
everywhere belongs to a point mass.

If \(n=3\), equation (*) is
\(\int q_\mu(t)q_\mu(2t)\,d\lambda=c(t)\).
Cauchy--Schwarz and the preceding squared-modulus identity give
\[
                              c(t)^2\leq c(t)c(2t).
\]
Continuity gives a neighborhood of zero where \(c\) is positive.
For any \(s\) in that neighborhood, repeated application at
\(s/2,s/4,\ldots\) gives \(c(s/2^k)\leq c(s)\).
Letting \(k\to\infty\) gives \(1\leq c(s)\leq1\).
Apply the elementary fact to the original pair \(Z_1,Z_2\):
its difference transform is \(c\), so the input law is degenerate,
a contradiction.

For \(n\geq4\), take real parts in (*) and use the modulus bound:
\[
 c(t)=\operatorname{Re}\int q_\mu(t)^{n-2}q_\mu(2t)\,d\lambda
       \leq\int |q_\mu(t)|^{n-2}|q_\mu(2t)|\,d\lambda
       \leq\int|q_\mu(t)|^2\,d\lambda=c(t).
\]
The nonnegative difference between the last two integrands has
integral zero. It follows that for almost every \(\mu\),
\[
                  q_\mu(t)\ne0\quad\Longrightarrow\quad
                                      |q_\mu(2t)|=1.
\]
Here division by \(|q_\mu(t)|^2\) is made only where it is positive;
the equality requires
\(|q_\mu(t)|^{n-4}|q_\mu(2t)|=1\), which gives the implication.
Remove the union of these null sets for rational \(t\).
For each remaining \(\mu\) and each rational \(s\), continuity at zero
gives \(q_\mu(s/2^k)\ne0\) for sufficiently large \(k\).
Iterating the implication gives \(|q_\mu(s)|=1\).
Continuity and density of the rationals extend this to every real \(s\).
The elementary fact shows that each such \(\mu\) is a point mass.
Consequently its product law makes \(Y_1=Y_2\) almost surely,
and integration over \(\lambda\) preserves that assertion.
But \(Y_1-Y_2=Z_1-Z_2\), so the original i.i.d. pair is equal
almost surely, again contradicting nondegeneracy.
This proves the full residual obstruction for every \(n\geq3\).

When the original variance is finite and positive, its simpler
numerical obstruction can also be calculated:
\[
 \operatorname{Var}(\overline Z)=\sigma^2/n,\quad
 \operatorname{Cov}(Z_i,\overline Z)=\sigma^2/n,\quad
 \operatorname{Cov}(Z_i,Z_j)=0\ (i\ne j).
\]
Thus \(EY_i=0\) and
\(\operatorname{Cov}(Y_i,Y_j)=-\sigma^2/n\) for distinct retained
coordinates. This is an additional special-case calculation,
not the hypothesis of the general proof.

For \(n=2\), only one residual is retained.
Any one-variable law \(\eta\) is the mixture of point masses
\(\eta(B)=\int\delta_y(B)\,\eta(dy)\); hence the claimed obstruction
does not apply to that one-coordinate vector.
If all inputs are equal to a constant almost surely, every residual
is zero and the law is itself a degenerate independent product.
These show why the dimension and nondegeneracy hypotheses of the
full residual obstruction cannot be omitted.
\end{proof}
\end{annexitem}

\begin{annexitem}{annex:p15.feller-approximation-boundaries}{thm:p15-feller-approximation-boundaries}{Explicit failures outside the approximation and positivity domains}
\begin{proof}
For the first law the signed mean is zero by symmetry, and
the second moment is \(n^2/n^3=1/n\).
The fourth moment is \(n^4/n^3=n\), proving the failure for the
unbounded continuous test.

For the second example define on \(\mathbb R\)
\[
                       u(x)=\sum_{m=1}^\infty
                                      \max(0,1-10m|x-m|).
\]
The summands are supported on
\([m-1/(10m),m+1/(10m)]\).
These intervals are disjoint and locally finite, so the sum is
continuous and takes values in \([0,1]\).
It satisfies \(u(n)=1\).
For \(n\geq2\), the points \(n\pm1/n\) lie outside the \(n\)th
support. Their distances to every other positive integer are at
least \(1-1/n\geq1/2\), exceeding every other support radius.
Thus \(u(n-1/n)=u(n+1/n)=0\).
The stated symmetric law has mean \(\theta\) and variance \(1/n^2\)
for every \(\theta\), but at \(\theta=n\) its expectation of \(u\)
is zero, while \(u(\theta)=1\).
Uniform convergence on the whole parameter line therefore fails.
The function is not uniformly continuous: its unit-height peaks
have widths tending to zero. The compact-parameter theorem is intact.

For \(u=-1\), every positive-order derivative is zero, but
the order-zero nonnegativity condition fails.
For \(u(x)=x-x^2\), a direct calculation gives
\(\Delta_h^2u(0)=-2\) whenever the required grid is in the interval;
in particular the source's \(k=2,n\geq3\) condition fails.
Finally the geometric series gives \((1-x)^{-1}=\sum_{r\geq0}x^r\)
for \(0\leq x<1\), with nonnegative coefficients whose sum is infinite.
It is not a probability generating function.
The zero function has all coefficients zero and satisfies all three
analytic conditions, but \(u(a)=0\) at every \(a>0\); a proof dividing
by such a value would not cover it. The subinterval coefficient proof
above includes it without that division.
\end{proof}
\end{annexitem}
